\BourbakiExercises{Exercises}

\begin{enumerate}
\def\labelenumi{\arabic{enumi})}
\tightlist
\item
  Let M be an A-module and N a direct factor submodule of M.
\end{enumerate}

\begin{enumerate}
\def\labelenumi{\alph{enumi})}
\item
  Prove that N is stable under the bicommutant \(A_{M}^{\prime\prime}\) of M.
\item
  Prove that every A-linear mapping from N to M is \(A_{M}^{\prime\prime}\) -linear.
\item
  Prove that restriction to N defines a ring homomorphism \(A_{M}^{\prime\prime} \to A_{N}^{\prime\prime}\) .
\item
  Suppose that M/N is a sum of submodules isomorphic to quotients of N. Prove that the canonical homomorphism \(A_{M}^{\prime\prime} \rightarrow A_{N}^{\prime\prime}\) is injective. If N is balanced, the same holds for M.
\end{enumerate}

\begin{enumerate}
\def\labelenumi{\arabic{enumi})}
\setcounter{enumi}{1}
\tightlist
\item
  Let K be a commutative field and A the subring of \(\mathbf{M}_{3}(\mathrm{K})\) consisting of the matrices of the form
\end{enumerate}

\[
\left( \begin{array}{c c c} a & 0 & 0 \\ b & c & 0 \\ 0 & 0 & a \end{array} \right)
\]

for \(a, b, c\) in K. Let M be the A-module \(\mathrm{K}^3\) ; it is the direct sum of the submodules \(\mathrm{N} = \mathrm{Ke}_1 + \mathrm{Ke}_2\) and \(\mathrm{P} = \mathrm{Ke}_3\) .

\begin{enumerate}
\def\labelenumi{\alph{enumi})}
\item
  Prove that M is balanced, that the canonical mapping \(A_{M}^{\prime\prime} \rightarrow A_{N}^{\prime\prime}\) (Exercise 1) is injective but not surjective, and that the canonical mapping \(A_{M}^{\prime\prime} \rightarrow A_{P}^{\prime\prime}\) is surjective but not injective.
\item
  Deduce an example of a module Q over a ring B and a direct factor submodule R such that the canonical mapping \(B_{Q}^{\prime\prime} \rightarrow B_{R}^{\prime\prime}\) is neither injective nor surjective (take \(B = A \times A\) and \(Q = M \times M\) ).
\end{enumerate}

\begin{enumerate}
\def\labelenumi{\arabic{enumi})}
\setcounter{enumi}{2}
\item
  Let B be a ring and I a two-sided ideal of B; let A be the subring of \(B \times B\) consisting of the pairs \((x, y)\) such that \(x - y \in I\) . Let \(B_{1}\) (resp. \(B_{2}\) ) be the additive group B endowed with an A-module structure deduced from the first (resp. second) projection. Prove that the A-modules \(B_{1}\) and \(B_{2}\) are balanced. Under what condition on I is the A-module \(B_{1} \oplus B_{2}\) balanced?
\item
  Let A and B be rings and P an (A, B)-bimodule.
\end{enumerate}

\begin{enumerate}
\def\labelenumi{\alph{enumi})}
\item
  Suppose that the B-module P is generating and that the A-module P is balanced. Show that for every balanced B-module M, the A-module \(P \otimes_{B} M\) is balanced.
\item
  Suppose that the A-module P is generating and that the mapping \(b \mapsto b_{P}\) from B to \(\operatorname{End}_{\mathrm{A}}(\mathrm{P})\) is bijective. Prove that for every generating B-module M, the A-module \(P \otimes_{B} M\) is generating.
\end{enumerate}

\begin{enumerate}
\def\labelenumi{\arabic{enumi})}
\setcounter{enumi}{4}
\item
  Let A be a ring, M an A-module, and N a submodule of M. Prove that if the A-module M/N is generating, the same holds for the A-module M. Use an example to prove that N can be generating without M being so.
\item
  Let \((\mathrm{A}_i)_{i\in \mathrm{I}}\) be a family of rings and A be its product. Prove that the A-module \(\bigoplus \mathrm{A}_i\) is projective and faithful. Is it always generating?
\item
  Let A be a ring and M an A-module. Prove that if M is generating, the right A-module \(\mathrm{M}^{*} = \mathrm{Hom}_{\mathrm{A}}(\mathrm{M}, \mathrm{A}_{s})\) is generating. Use an example to prove that \(M^{*}\) can be generating without M being so (cf.~Exercise 6).
\item
  Let A be a principal ideal domain and M an A-module. Prove that the following properties are equivalent:
\end{enumerate}

\begin{enumerate}
\def\labelenumi{(\roman{enumi})}
\item
  The A-module M is generating.
\item
  The dual of M is nonzero.
\item
  The module M contains a direct factor submodule isomorphic to \(A_{s}\) .
\end{enumerate}

\begin{enumerate}
\def\labelenumi{\arabic{enumi})}
\setcounter{enumi}{8}
\tightlist
\item
  Let \(D\) be a field and \(V\) a right vector space over \(D\) . We view \(V\) as a left module over the ring \(A = \operatorname{End}_D(V)\) .
\end{enumerate}

\begin{enumerate}
\def\labelenumi{\alph{enumi})}
\item
  Prove that the A-module V is projective, finitely generated, faithful, and balanced. Prove that the (D,A)-bimodules \(\mathrm{Hom}_{\mathrm{A}}(\mathrm{V},\mathrm{A}_{s})\) and \(\mathrm{V}^{*}=\mathrm{Hom}_{\mathrm{D}}(\mathrm{V},\mathrm{D}_{d})\) are canonically isomorphic.
\item
  Prove that the A-module V is generating if and only if the dimension of the D-vector space V is finite.
\item
  What is the trace ideal of the A-module V?
\end{enumerate}

\begin{enumerate}
\def\labelenumi{\arabic{enumi})}
\setcounter{enumi}{9}
\tightlist
\item
  Let M be an A-module. Prove that the following properties are equivalent:
\end{enumerate}

\begin{enumerate}
\def\labelenumi{(\roman{enumi})}
\item
  The A-module M is finitely generated.
\item
  For every generating A-module P, there exist a natural number n and a surjective homomorphism f from \(P^{n}\) to M.
\end{enumerate}

\begin{enumerate}
\def\labelenumi{\arabic{enumi})}
\setcounter{enumi}{10}
\tightlist
\item
  Let \(\mathrm{P}\) be a finitely generated projective A-module; denote its trace ideal by \(\tau(\mathrm{P})\) .\\
\end{enumerate}

\begin{enumerate}
\def\labelenumi{\alph{enumi})}
\item
  Prove that we have \(\tau(\mathrm{P})^2 = \tau(\mathrm{P})\) and \(\tau(\mathrm{P})\mathrm{P} = \mathrm{P}\) .
\item
  Prove that the canonical mapping \(\tau(\mathrm{P})\otimes_{\mathrm{A}}\mathrm{P}\to\mathrm{P}\) is an isomorphism.
\end{enumerate}

¶ 12) Let A be a ring, and let P be a finitely generated A-module. Prove that the following properties are equivalent:

\begin{enumerate}
\def\labelenumi{(\roman{enumi})}
\item
  The A-module P is projective and generating.
\item
  The A-modules \(\mathrm{P}^{(\mathbf{N})}\) and \(\mathrm{A}_s^{(\mathbf{N})}\) are isomorphic.
\end{enumerate}

(To establish the implication \((\mathrm{i})\Rightarrow(\mathrm{ii})\) , use induction to construct strictly increasing sequences \((n_{k})\) and \((m_{k})\) of integers and surjective homomorphisms of A-modules \(f_{k}:A_{s}^{n_{k}}\to P^{m_{k}}\) and \(g_{k}:P^{m_{k+1}}\to A_{s}^{n_{k}}\) such that \(g_{k-1}\circ f_{k}\) and \(f_{k}\circ g_{k}\) are the canonical projections.)

\begin{enumerate}
\def\labelenumi{\arabic{enumi})}
\setcounter{enumi}{12}
\item
  Let k be a ring and n an integer \(\geqslant 2\) ; denote by A the subring of \(\mathbf{M}_{n}(k)\) consisting of the upper-triangular matrices and by M the right A-module \(k^{n}\) . Prove that M is projective and finitely generated but that the ring of homotheties of M is not dense in its bicommutant (cf.~VIII, p.~88, Theorem 3 and VIII, p.~88, Remark). In particular, M is not balanced.
\item
  What is the bicommutant of the Z-module Q? Prove that Q is a simple module over its bicommutant but that the ring of homotheties of the Z-module Q is not dense in its bicommutant.
\end{enumerate}

¶ 15) Let M be a torsion module over a principal ideal domain. Prove that the ring of homotheties of M is dense in its bicommutant.

(Reduce to the case when M is \(\pi\) -primary for an irreducible element \(\pi\) of A. Then, prove that every finite subset S of M is contained in a direct factor submodule that is the direct sum of a finitely generated module and a finite number of indecomposable divisible modules. Use induction on the cardinal of S with the help of Exercises 3 of VII, §2, p.~54 and 8 of VII, §2, p.~56.)

\begin{enumerate}
\def\labelenumi{\arabic{enumi})}
\setcounter{enumi}{15}
\tightlist
\item
  Let D be a field, V a vector space of dimension \(\geqslant 2\) over D, and A a subring of \(\operatorname{End}_{\mathrm{D}}(\mathrm{V})\) .
\end{enumerate}

\begin{enumerate}
\def\labelenumi{\alph{enumi})}
\item
  Suppose that for every system of four elements \((x, x', y, y')\) , where x and \(x'\) are linearly independent, there exists an element u of A such that \(u(x) = y\) and \(u(x') = y'\) . Prove that the commutant of the A-module V is D (identified with the ring of homotheties of V) and that the ring of homotheties of this module is dense in its bicommutant.
\item
  Give an example of a ring \(A \subset \operatorname{End}_{\mathrm{D}}(\mathrm{V})\) such that V is a simple A-module but that its commutant is distinct from D (consider a field E containing D, a vector space V over E, and the ring \(A = \operatorname{End}_{\mathrm{E}}(\mathrm{V})\) ).
\end{enumerate}

\begin{enumerate}
\def\labelenumi{\arabic{enumi})}
\setcounter{enumi}{16}
\tightlist
\item
  Let K be an algebraically closed commutative field, V a finite-dimensional vector space over K, and G a submonoid of \(\operatorname{Aut}(V)\) such that the K{[}G{]}-module V is simple. Prove that if the set of scalars \(\operatorname{Tr}(g)\) , for \(g \in G\) , is finite, then G is finite (deduce from Burnside's theorem that G contains a basis \((g_{i})_{i \in I}\) of \(\operatorname{End}(V)\) , and consider the mapping \(g \mapsto (\operatorname{Tr}(gg_{i}))\) from G to \(K^{I}\) ).
\end{enumerate}

¶ 18) A group G is called locally finite if every subgroup of G that admits a finite generating family is finite.

\begin{enumerate}
\def\labelenumi{\alph{enumi})}
\item
  Let G be a group and H a normal subgroup of G. Prove that if H and G/H are locally finite, the same holds for G.
\item
  Deduce from a) that a solvable group whose elements are all of finite order is locally finite.
\item
  Let K be a commutative field, V a finite-dimensional vector space over K, and G a subgroup of Aut(V) that admits a finite generating family consisting of elements of finite order. Prove that the order of the elements of G is bounded.
\end{enumerate}

(Reduce to the case when K is a purely transcendental extension of a prime field P by considering the subfield K generated by the coefficients of a finite generating family of G. Then, prove that for every element s of \(\mathbf{GL}_{n}(\mathrm{K})\) of finite order, we have \(\varphi(s) \leqslant n\) if P = Q and \(s \leqslant p^{n} - 1\) if \(P = F_{p}\) .)

\begin{enumerate}
\def\labelenumi{\alph{enumi})}
\setcounter{enumi}{3}
\tightlist
\item
  Prove that every subgroup H of Aut(V) whose elements are all of finite order is locally finite. \(^{(1)}\)
\end{enumerate}

(Reduce to the case when K is algebraically closed, and use induction on \(\dim(V)\) . Let G be a subgroup of H that admits a finite generating family. If V is a simple K{[}G{]}-module, apply Exercise 17; if V contains a nontrivial subspace W that is stable under G, consider the image of G in \(\operatorname{Aut}(V) \times \operatorname{Aut}(V/W)\) , and apply a) and b)).

¶ 19) Let A be a ring such that the A-module \(A_{s}\) is the direct sum of a finite family of left ideals \((\mathfrak{I}_{i})_{1\leqslant i\leqslant n}\) that are isomorphic as A-modules.

\begin{enumerate}
\def\labelenumi{\alph{enumi})}
\item
  Prove that in A, there exists a family of elements \(e_{i,j}\) ( \(1 \leqslant i \leqslant n\) , \(1 \leqslant j \leqslant n\) ) such that \(e_{i,j}e_{h,k} = \delta_{j,h}e_{i,k}\) , where \(\delta_{j,h}\) denotes the Kronecker delta function, and \(I_i = Ae_{i,i}\) (cf.~VIII, p.~39, Exercise 5 and I, §8, p.~173, Exercise 12). Conversely, if in A, there exists a family of \(n^2\) elements \(e_{i,j}\) such that \(e_{i,j}e_{h,k} = \delta_{j,h}e_{i,k}\) and \(1 = \sum_{i=1}^{n} e_{i,i}\) , then the left ideals \(Ae_{i,i}\) are isomorphic as A-modules (cf.~VIII, p.~39, Exercise 5). Moreover, if B is the subring of A consisting of the elements that commute with all the \(e_{i,j}\) , then A is isomorphic to the matrix ring \(\mathbf{M}_n(\mathbf{B})\) and B is isomorphic to the opposite ring of the commutant of each of the A-modules \(Ae_{i,i}\) (use VIII, p.~78 and VIII, p.~65, Corollary).
\item
  Let M be an A-module. Prove that the \(N_{i} = e_{i,i}M\) are B-modules and that M, viewed as a B-module, is the direct sum of the \(N_{i}\) ; moreover, the B-modules \(N_{i}\) are pairwise isomorphic (consider the mappings \(x \mapsto e_{i,1}x\) and \(y \mapsto e_{1,i}y\) ), and the annihilator of the \(N_{i}\) is the intersection of B and the annihilator of M in A. Conversely, for every B-module N, define an A-module structure on the direct sum M of n B-modules isomorphic to N, so that \(e_{1,1}M\) is isomorphic to N.
\item
  To each B-submodule P of \(N_{1}\) , assign the A-submodule \(\sum_{i=1}^{n} e_{i,1}P\) of M. Prove that this defines a bijective mapping from the set of B-submodules of \(N_{1}\) to the set of A-submodules of M that is strictly increasing for the inclusion relation.
\end{enumerate}

